%=====================================================================
\section{Two Annotation Tasks as Procedures}
\label{sec:tasks}

We test the method on two annotation schemes that differ in what the output is and how much of it is structure. In both, the seed procedure was written by hand from the published guideline, and the same driver, reflector and decision model learn it.

\paragraph{Verbal multiword expressions (PARSEME).}
The PARSEME shared tasks annotate verbal multiword expressions (VMWEs): sets of tokens, possibly discontinuous, with a category such as a verbal idiom (VID), a light-verb construction (LVC.full, LVC.cause), a verb--particle construction (VPC.full, VPC.semi), an inherently adpositional verb (IAV), or, for Chinese, a multi-verb construction (MVC) \citep{ramisch2018parseme,ramisch2020parseme}. In Chinese a VMWE is often a single token: a verb with its directional or resultative complement written as one word (\zht{推出}, \zht{放入}) or a four-character idiom. We use Chinese from edition 1.2 and English from edition 1.1, the last edition with English. The input gives each token's lemma, part of speech and dependency relation; these are part of the data, and the procedure reads them like any other input.

The seed procedure has two files. \texttt{candidates.py} proposes candidates by rule: in Chinese, a verb token ending in a complement (split into stem and complement), a verb and a following complement token, a light verb and its object, and a four-character token; in English, a verb and its particle, a light verb and its object noun, a verb and the preposition after it, and \emph{make/let/get/break} with a bare complement. \texttt{decide.py} asks one choice question per candidate, whose options are the categories the candidate's kind allows plus \emph{none}, each described in one line from the guideline; a candidate is kept when its best option is a category with probability $\geq 0.5$, and overlaps are resolved by probability. The seed's candidates cover about three quarters of the gold VMWEs of the training data. All questions of a sentence go to Jev in one request.

\paragraph{Uniform Meaning Representation (UMR).}
UMR \citep{bonn2023umr,bonn2025umr2} represents each sentence as a rooted graph of concepts (PropBank senses for events, words or entity types for objects, abstract predicates such as \texttt{identity-91}) linked by roles, with attributes such as aspect and polarity, plus document-level temporal, modal and coreference relations. The output is mostly structure. The procedure splits a sentence top-down. A routing skill decides whether a span is coordinated clauses, a clause, a noun phrase, a named entity or a leaf; a clause skill asks for the predicate, its sense, aspect, modal strength and negation, and the role of every word in the clause; noun-phrase and entity skills ask for the head, the shape and the modifiers' relations. Unexpanded parts of a fragment become placeholders that go back to routing, rules cover words that no child took, and a reentrancy pass merges coreferent nodes. A sentence takes 11--49 requests to Jev, depending on the procedure's version. We start from the closed procedure of the earlier version of this work and from its recent revision, whose questions and options were moved into data files and whose relation inventories were completed (\S\ref{sec:analysis}). We score sentence graphs with AnCast++ \citep{sun2024ancast}; document-level relations are outside this study.

\paragraph{Data.}
For PARSEME we use the official files: Chinese train (35,326 sentences), dev and test (3,462 sentences, 786 VMWEs, 38\% of them unseen in train); English train (3,471) and test (3,965 sentences, 501 VMWEs). Reflection reads 5-sentence batches from the training file, three sentences with a VMWE and two without, so that every batch has something to learn from; the development set is 150 sentences in their natural proportion (from the official dev file for Chinese, held out of train for English, which has no dev file). For UMR we use the 31 Chinese gold documents and the native UMR 2.0 English annotation, excluding the English sentences that were converted only in part from other banks. Chinese: batches from documents 0004--0014 (taken round-robin over the documents), a development set of 27 sentences, and a test set of eight documents (138 sentences) that earlier analyses of this project had looked at, so the Chinese UMR comparison is retrospective. English: batches from the guideline's example sentences and a narrative document, 25 development sentences, and a test set of five news documents and a narrative passage (143 sentences). Twelve rounds read 60 training sentences.
